\section{Certificate layers A6--A7d}
\label{sec:layers}

This section records the certificate layers used by Theorem~\ref{thm:wrapper}.
The layers are package-local and source-backed; their role is to discharge the
routed pair obligations, not to prove a broader mechanism theorem.

\begin{certprop}[A6: B05 common-edge projection]
The B05 common-edge projection layer is symbolically closed on the one-parameter
ray for the relevant records.  A6 packages the A3/A4 and A5 algebraic
certificates: A3/A4 proves full-domain raw gap, normalized gap, and axis
positivity on the open ray, while A5 splits at \(t=1\) and proves post-switch
support signatures and projection gaps on \(1<t<7/4\), which contains
\(1<t\le\sqrt3\).  The A7d wrapper uses four B05 pair routes.
\end{certprop}

\begin{proof}[Certification]
The source is
\path{certified/source_docs/S4_CL5_A6_ONE_PARAMETER_RAY_CLOSURE_PACKAGE.md}.
It records seven symbolic B05 closures in the source package and the one-
parameter B05 closure chain used by the wrapper.  The public wrapper consumes
the B05 role through \path{certified/a7d_one_parameter_theorem_wrapper_manifest.json}.
\end{proof}

\begin{certprop}[A7a: B06/B07 shared-face residual]
For the two one-parameter shared-face residual targets,
\[
\treeA{}:P_2-P_3,
\qquad
\treeB{}:P_0-P_2,
\]
the residual shared-face formula is positive on the open ray.  In half-angle
coordinate the certified rational form is
\[
\frac{t^3}{(1+t^2)^2},
\qquad 0<t<2.
\]
In particular it is positive on \(0<t\le\sqrt3\).
\end{certprop}

\begin{proof}
The trigonometric formula from the source layer is
\[
\frac{(1-\cos\theta)\sin\theta}{4}
  =\sin^3(\theta/2)\cos(\theta/2).
\]
With \(t=\tan(\theta/2)\), this becomes
\[
\frac{t^3}{(1+t^2)^2}.
\]
The numerator and denominator have no open roots on \(0<t<2\); at \(t=1\) both
are positive.  Therefore the expression is positive on the certified open
superset, hence on the audited ray domain.  The Sturm/root-count certificate is
recorded in
\path{certified/source_docs/S4_CL5_A7A_SHARED_FACE_RESIDUAL_STURM_CERTIFICATE.md}.
\end{proof}

\begin{certprop}[A7b: B03 route vacuity]
There are no route-clean ordinary non-contact B03 obligations on the audited
one-parameter representative rays.  Across the two representatives, the twelve
pair obligations route as six selected-hinge contacts, four residual shared-edge
common-edge projections, and two residual shared-face targets.
\end{certprop}

\begin{proof}[Certification]
This is the A7b route-vacuity certificate in
\path{certified/source_docs/S4_CL5_A7B_B03_RAY_VACUITY_CERTIFICATE.md}.  It
records \texttt{one-parameter pair records: 12}, \texttt{B03 route-clean pairs:
0}, and \texttt{B03 Sturm obligations: 0}.
\end{proof}

\begin{certprop}[A7c: selected-hinge signed orientation]
For the six selected B04 rows, A7c certifies the sign of the selected hinge
angle on the open ray.  The signed orientation has the same sign as \(s_h t\)
on \(0<t<2\), with ray-sign counts \(+1:4\) and \(-1:2\).
\end{certprop}

\begin{proof}
Lemma~\ref{lem:kinematics} fixes the signed-ray convention
\(d_h(\theta)=s_h\theta\).  In half-angle coordinates, the relevant signed
orientation is sign-equivalent to \(s_h t\).  The only root is at \(t=0\),
which is excluded from the open ray.  The sample point \(t=1\) fixes the sign.
The six rows are recorded in
\path{certified/a7c_selected_hinge_contact_side_certificate_manifest.json} and
\path{certified/source_docs/S4_CL5_A7C_SELECTED_HINGE_CONTACT_SIDE_CERTIFICATE.md}.
\end{proof}

A7c alone is not the B04 contact-side proof.  It is the signed-orientation
source layer consumed by the bridge in Section~\ref{sec:b04}.
